% ============================================================================
% CAPÍTULO 1: INTRODUCCIÓN
% ============================================================================
% 
% INSTRUCCIONES GENERALES:
% - Este capítulo presenta una visión general de la investigación
% - Debe ser claro, conciso y motivar al lector
% - Extensión típica: 10-15 páginas
% - Reemplaza TODO el texto entre [corchetes]
% ============================================================================

\chapter{Introducción}
\label{cap:introduccion}

% ----------------------------------------------------------------------------
% CONTENIDO REQUERIDO PARA INTRODUCCIÓN:
% ----------------------------------------------------------------------------
% Este capítulo debe incluir una visión general de:
% - El contexto y relevancia del tema de investigación
% - La problemática que abordas
% - Una breve descripción de lo que harás
% - La estructura del documento
%
% NOTA IMPORTANTE: La Delimitación del Objeto de Estudio (planteamiento del
%                  problema, pregunta de investigación, hipótesis, objetivos)
%                  va en el CAPÍTULO 3, NO aquí.
% ----------------------------------------------------------------------------

\noindent El Comportamiento Sedentario (CS) representa actualmente uno de los desafíos más críticos de salud pública del siglo XXI, vinculado estrechamente con la transición epidemiológica, la urbanización acelerada y el cambio en los patrones de movilidad y trabajo. Las largas jornadas en posición sentada, el predominio de actividades digitales y la disminución de la actividad física incidental han generado un aumento sostenido en la prevalencia de Enfermedades Crónicas No Transmisibles (ENT), como la diabetes tipo 2, las enfermedades cardiovasculares y la depresión \cite{Bull2020}. Este panorama, descrito por \cite{WHO2020}, refleja la urgencia de desarrollar herramientas que permitan medir, interpretar y reducir los efectos del sedentarismo en la Calidad de Vida Relacionada con la Salud (CVRS) \cite{Guthold2020}.

El estudio del CS tradicionalmente utiliza instrumentos de autoinforme, como el Cuestionario Mundial sobre Actividad Física (\textit{GPAQ}, por sus siglas en inglés) o el Cuestionario Internacional de Actividad Física (\textit{IPAQ}, por sus siglas en inglés), que aunque válidos y ampliamente implementados, presentan limitaciones inherentes a la subjetividad y la falta de resolución temporal \cite{Romero2022}. En la última década, el desarrollo de dispositivos portátiles con sensores biométricos, que incluyen acelerómetros, fotopletismografía (\textit{PPG}) y sensores de frecuencia cardiaca, ha permitido capturar datos fisiológicos y de actividad con alta precisión; lo cual abre la posibilidad de análisis objetivos, continuos y personalizados \cite{Henriksen2018,White2019,Strain2020}. Este desarrollo tecnológico sitúa al monitoreo digital de la salud como una pieza central en el diseño de intervenciones preventivas basadas en evidencia.

En este contexto, la Inteligencia Artificial (IA) y, particularmente, la lógica difusa, ofrecen un marco teórico y computacional idóneo para modelar fenómenos complejos caracterizados por relaciones no lineales entre variables. La lógica difusa admite grados intermedios de pertenencia y permite representar razonamientos humanos en términos lingüísticos, lo cual otorga interpretabilidad y flexibilidad a los modelos biomédicos \cite{Zadeh1965,Ross2010,Kaur2022}. Su integración en sistemas de inferencia \textit{Mamdani} posibilita traducir datos biométricos en reglas semánticamente comprensibles y fortalece la capacidad explicativa de los algoritmos.

El presente proyecto desarrolló un modelo de evaluación del Comportamiento Sedentario basado en lógica difusa e información biométrica proveniente del \textit{Apple Watch}, que integra variables de actividad física (minutos en movimiento, pasos diarios, gasto calórico) con indicadores autonómicos derivados de la variabilidad de la frecuencia cardiaca (\textit{HRV-SDNN}). Este enfoque caracterizó la relación entre patrones de sedentarismo y percepción de CVRS, capturando la dinámica diaria del equilibrio entre esfuerzo, recuperación y bienestar. A través de un sistema de inferencia difusa con reglas adaptativas, se modeló la interacción entre la activación fisiológica y los hábitos conductuales, lo cual ofreció una aproximación más precisa y contextualizada del impacto del sedentarismo sobre la salud.

Desde una perspectiva metodológica, el modelo difuso se nutre de datos reales obtenidos de Apple \textit{HealthKit}, los cuales son procesados mediante un protocolo reproducible que garantiza la integridad y trazabilidad de la información. Se construye una base de control difuso con funciones de pertenencia triangulares y trapezoidales para las variables de entrada, y un conjunto de reglas lingüísticas. El sistema se valida mediante concordancia con una verdad operativa derivada de \textit{clustering} no supervisado, y utiliza métricas de rendimiento (\textit{F1-Score}, Exactitud, Coeficiente de Matthews) y validación cruzada \textit{Leave-One-User-Out} (\textit{LOUO}) para evaluar generalización inter-sujeto, es decir, la capacidad del modelo para clasificar correctamente el comportamiento sedentario en usuarios no incluidos en el entrenamiento. Adicionalmente, se realizó una validación convergente exploratoria con el cuestionario \textit{SF-36} en un subconjunto de participantes (N=8), la cual reveló correlación significativa únicamente con la dimensión salud mental.

Este modelo contribuye a la comprensión del sedentarismo como fenómeno fisiológico y conductual. Integra la dimensión autonómica del organismo en el análisis de la actividad física. El sistema es explicable y flexible. Se adapta a distintas poblaciones y contextos de uso. La investigación aporta un marco epistemológico para diseñar estrategias de promoción de la salud personalizadas y sostenibles. Estas estrategias se alinean con la agenda 2030 para el desarrollo sostenible, específicamente con la meta 3.4 de los Objetivos de Desarrollo Sostenible (ODS), que busca reducir en un tercio la mortalidad prematura por ENT para 2030, y con el plan mundial sobre actividad física 2018–2030 \cite{WHO2018}.

La urgencia de abordar el CS se evidencia en datos globales alarmantes. En 2019, las ENT provocaron cinco millones de muertes a nivel global, y en la región de las Américas representan casi cuatro de cada cinco muertes anuales \cite{StrategicPlanPAHO2020}. En México, el CS en adultos alcanza el 40\% \cite{Campos2023ObesityMexico}. Los costos del CS se estiman en más de 300 000 millones de dólares anuales a nivel mundial, debido a costos directos de atención médica y pérdida de productividad económica \cite{Okunogbe2021Economic}. Este proyecto articula tres dimensiones convergentes: la necesidad epidemiológica de reducir el CS, el potencial tecnológico de los dispositivos portátiles para capturar datos de alta resolución y la capacidad analítica de la lógica difusa para transformar esos datos en conocimiento clínicamente significativo. Con ello, se establece un precedente metodológico para futuras investigaciones en salud digital, inteligencia artificial y promoción de la salud basada en evidencia.

% ----------------------------------------------------------------------------
% SECCIÓN: Organización del Documento
% ----------------------------------------------------------------------------
% [Eliminado: Organización del Documento]

% ============================================================================
% NOTAS FINALES:
% ============================================================================
% - Recuerda citar fuentes relevantes usando \cite{clave_referencia}
% - Mantén un tono académico pero accesible
% - Evita entrar en detalles metodológicos (eso va en Cap. 5)
% - No incluyas objetivos ni hipótesis aquí (eso va en Cap. 3)
% ============================================================================
